\documentclass[12pt,a4paper]{article}
\usepackage{amsmath,amssymb,amsthm}
\usepackage{booktabs}
\usepackage{hyperref}
\usepackage{enumitem}
\usepackage{geometry}
\geometry{margin=2.5cm}

\newtheorem{theorem}{Theorem}[section]
\newtheorem{corollary}[theorem]{Corollary}
\theoremstyle{definition}
\newtheorem{definition}[theorem]{Definition}
\theoremstyle{remark}
\newtheorem{remark}[theorem]{Remark}

\title{The Boltzmann Constant from Recognition Science:\\
$k_R = \ln\varphi$ as the Ledger Bit-Cost}

\author{Jonathan Washburn\\
Recognition Science Research Institute\\
Austin, Texas\\
\texttt{washburn.jonathan@gmail.com}}

\date{March 2026}

\begin{document}
\maketitle

\begin{abstract}
We derive the Recognition Science analog of the Boltzmann constant
from the cost functional $J(x) = \tfrac{1}{2}(x + x^{-1}) - 1$.
The elementary bit-cost---the $J$-cost of one binary distinction on
the $\varphi$-ladder---is $J_{\mathrm{bit}} = \ln\varphi$.  The RS
thermal constant $k_R = \ln\varphi \approx 0.481$ relates temperature
to energy: $E = k_R\,T$.  We prove $k_R > 0$, $k_R < 1/2$, and
$0.47 < k_R < 0.49$.  The identification $k_R = J_{\mathrm{bit}}$
shows that the Boltzmann constant is not an independent thermodynamic
parameter but the cost of one bit of information on the recognition
ledger.  The SI Boltzmann constant
$k_B = 1.381 \times 10^{-23}\;\mathrm{J/K}$ is recovered through the
calibration seam (a temperature-scale anchor maps RS units to Kelvin).
\end{abstract}

\section{Introduction}

The Boltzmann constant $k_B$ bridges the microscopic and macroscopic
descriptions of matter, converting between temperature and energy.  In
SI, it is defined to be exactly
$1.380\,649 \times 10^{-23}\;\mathrm{J/K}$.

Recognition Science~\cite{RS_Axioms} derives a natural thermal
constant $k_R$ from the $\varphi$-ladder structure.  The key insight:
temperature measures the number of ledger bits excited, and the cost
per bit is $\ln\varphi$.

\section{Recognition Science Framework}
\label{sec:framework}

Recognition Science derives all physics from a single primitive,
the \emph{recognition cost functional} $J(x)$, uniquely determined
by three axioms.

\begin{definition}[RS Cost Functional]
For $x > 0$: $J(x) = \tfrac{1}{2}(x + x^{-1}) - 1$.
\end{definition}

\noindent\textbf{Axiom A1 (Normalization):} $J(1) = 0$.\\
\textbf{Axiom A2 (Recognition Composition Law, RCL):}
$J(xy) + J(x/y) = 2J(x)J(y) + 2J(x) + 2J(y)$ for all $x, y > 0$.\\
\textbf{Axiom A3 (Calibration):} $J''_{\log}(0) = 1$, where
$J_{\log}(t) := J(e^t)$.

\begin{theorem}[Cost Uniqueness]
\label{thm:unique}
The unique continuous solution to \textup{(A1)--(A3)} is
$J(x) = \tfrac{1}{2}(x + x^{-1}) - 1$.
\end{theorem}

\begin{proof}[Proof sketch]
Setting $f(t) = J_{\log}(t) + 1$ and rewriting \textup{(A2)} gives the
d'Alembert equation $f(s+t) + f(s-t) = 2f(s)f(t)$.  By the Acz\'el
classification theorem~\cite{Aczel1966}, the unique continuous solution
with $f(0) = 1$ and $f''(0) = 1$ is $f(t) = \cosh t$.
\end{proof}

\noindent The golden ratio $\varphi = (1+\sqrt{5})/2$ is forced by the RCL
self-similarity equation.  The $\varphi$-ladder has multiplicative spacing
$\varphi$: one rung step corresponds to multiplying the recognition ratio
by $\varphi$.  The information content per rung step in natural logarithmic
units is $\ln\varphi \approx 0.481$.  This is the fundamental bit-cost of
the ledger.

\section{The Bit-Cost}

\begin{definition}[Ledger Bit-Cost]
The elementary bit-cost on the $\varphi$-ladder is
\begin{equation}
  J_{\mathrm{bit}} = \ln\varphi,
\end{equation}
where $\varphi = (1+\sqrt{5})/2$.
\end{definition}

\begin{remark}
This arises because the $\varphi$-ladder has multiplicative spacing
$\varphi$: moving one rung corresponds to multiplying the ratio by
$\varphi$.  The information content (in natural logarithmic units) of
one rung step is $\ln\varphi$.
\end{remark}

\section{The RS Thermal Constant}

\begin{definition}[RS Boltzmann Analog]
\begin{equation}
  \boxed{k_R = \ln\varphi \approx 0.481.}
\end{equation}
\end{definition}

\begin{remark}[$k_R$ Properties]
\begin{enumerate}
  \item $k_R > 0$ (since $\varphi > 1$, $\ln\varphi > 0$).
  \item $k_R < 1/2$ (since $\varphi < e^{1/2} \approx 1.649$).
  \item $0.476 < k_R < 0.483$ (from $\ln 1.61 < \ln\varphi < \ln 1.62$).
  \item $k_R = J_{\mathrm{bit}}$ (by definition).
\end{enumerate}
The exact value is $\ln\varphi = \ln((1+\sqrt{5})/2) \approx 0.48121$.
\end{remark}

\begin{remark}[Thermal Energy at Unit Temperature]
At temperature $T = 1$ (RS-native units), $E = k_R \cdot 1 = k_R = \ln\varphi$.
More generally, $E = k_R T$: temperature $T$ in RS-native units equals
the number of $\varphi$-ladder rungs thermally excited, and each costs
$k_R = \ln\varphi$.
\end{remark}

\section{Physical Interpretation}

The RS identification $k_R = \ln\varphi = J_{\mathrm{bit}}$ has a
clear physical meaning:
\begin{itemize}
  \item Temperature measures the average number of $\varphi$-ladder
  rungs that are thermally excited.
  \item Each rung carries information content $\ln\varphi$.
  \item The thermal energy is the total bit-cost: $E = k_R\,T =
  (\ln\varphi)\,T$.
\end{itemize}

This dissolves the ``mystery'' of the Boltzmann constant: it is not
a conversion factor between human-defined temperature and energy
units, but the fundamental cost of one bit of ledger information.

\section{Entropy and the $\varphi$-Ladder}

In RS, entropy counts the number of $\varphi$-rungs accessible to a
system.  For a system with $W$ microstates, the RS entropy is
\begin{equation}
  S = k_R \cdot \ln W = (\ln\varphi)\,\ln W.
\end{equation}
Each microstate corresponds to a configuration on the
$\varphi$-ladder; the logarithm $\ln W$ measures how many
distinguishable rungs the system can occupy.  Multiplying by $k_R =
\ln\varphi$ converts this count into the total ledger cost: the
entropy is the $J$-cost of the system's accessible configurations.

Equivalently, $S$ is the number of $\varphi$-rungs spanned by the
system's accessible phase space.  A system with $W$ equiprobable
microstates has $\ln W$ bits of information; at $\ln\varphi$ cost
per bit, the entropy $S = k_R \ln W$ is the total cost of that
information on the ledger.  High entropy means many rungs
accessible; low entropy means few.  The $\varphi$-ladder structure
thus provides a geometric interpretation of Boltzmann's formula
$S = k \ln W$.

\section{Connection to Information Theory}

The bit-cost $\ln\varphi$ bridges Shannon entropy and physical
entropy.  Shannon entropy (in nats) is
$H = -\sum_i p_i \ln p_i$.  For equiprobable outcomes,
$H = \ln W$.  In RS, each unit of Shannon entropy (one nat) costs
$\ln\varphi$ on the ledger.  Hence physical entropy
$S = k_R \cdot H = (\ln\varphi)\,H$.

The entropy per $\varphi$-rung is exactly one natural bit in
base-$\varphi$: $\log_\varphi(e^{\ln\varphi}) = 1$.  That is, one
step on the $\varphi$-ladder carries precisely one ``natural bit''
when information is measured in base-$\varphi$ units.  The RS
Boltzmann constant $k_R = \ln\varphi$ is thus the conversion factor
between information-theoretic entropy (nats) and thermodynamic
entropy (ledger cost).  Shannon and Boltzmann entropy are unified:
both count information; $k_R$ assigns the cost.

\section{Thermodynamic Laws}

The first and second laws of thermodynamics emerge from
$J$-cost minimization on the ledger.

\paragraph{First law (energy conservation).}
The ledger is balanced: inputs and outputs sum to zero.  Energy
cannot be created or destroyed because the $J$-cost of every
transaction is accounted for.  The first law is ledger balance---
$\Delta E = Q - W$ reflects the conservation of $J$-cost across
heat and work exchanges.

\paragraph{Second law (entropy increase).}
The $J$-cost functional $J(x) = \frac{1}{2}(x + x^{-1}) - 1 \geq 0$
has a unique minimum at $x = 1$.  Isolated systems undergo recognition
dynamics that minimize $J$-cost: the sequence of recognition events at
each tick drives the system toward configurations of lower $J$-cost.
In the thermodynamic limit, the number of accessible microstates
$W$ increases (by ergodicity of the recognition process) while the
$J$-cost per event decreases.  Since entropy $S = k_R \ln W$ is an
increasing function of $W$, the direction of increasing entropy ($W$
increasing) is identified with the direction of decreasing $J$-cost
(the system approaching the zero-defect equilibrium $x = 1$).
The approach to thermodynamic equilibrium (maximum $S$, minimum
$J$-cost) is thus structurally identified with RS cost-minimization
dynamics.  (The identification of $W \to \infty$ with $x \to 1$ is
the structural claim that needs derivation from the RS Hamiltonian,
as noted in the Remark below.)

\begin{remark}[Status of second law derivation]
The argument above is structural: it depends on the ledger being
ergodic in the thermodynamic limit, which is a plausible assumption
but not yet a theorem from the RS axioms.  A fully rigorous derivation
of the second law from RS (analogous to Boltzmann's $H$-theorem from
kinetic theory) remains an open problem.
\end{remark}

\section{Comparison and Predictions}

\begin{table}[h]
\centering
\begin{tabular}{@{}lll@{}}
\toprule
Quantity & Conventional & Recognition Science \\
\midrule
Thermal constant & $k_B$ (SI, measured) & $k_R = \ln\varphi \approx 0.481$ (derived) \\
Bit-cost & (implicit in stat.\ mech.) & $J_{\mathrm{bit}} = \ln\varphi$ \\
Entropy formula & $S = k_B \ln W$ & $S = k_R \ln W = (\ln\varphi)\ln W$ \\
Temperature--energy & $E = k_B T$ & $E = k_R T$ \\
Information basis & base-$e$ (nats) & base-$\varphi$ \\
SI value of constant & $1.381 \times 10^{-23}$ J/K & via calibration seam \\
\bottomrule
\end{tabular}
\caption{RS vs.\ conventional thermodynamics.  The RS thermal
constant $k_R = \ln\varphi$ is the \emph{structural} value;
the SI Boltzmann constant $k_B$ is recovered by anchoring the
RS temperature scale to the Kelvin through one calibration point.}
\end{table}

\begin{remark}[Relation to SI $k_B$]
The RS derivation predicts the \emph{structure} of the thermal
constant: it is $\ln\varphi$, the log-spacing of the
$\varphi$-ladder.  The SI value $k_B = 1.381 \times 10^{-23}$ J/K
involves the Kelvin (historically defined by water's triple point),
which introduces a human-defined anchor.  The ratio
$k_R / k_B^{\mathrm{natural}}$ (where $k_B^{\mathrm{natural}}$
is $k_B$ in Planck units $= k_B / (E_P/T_P)$) is the RS-native
prediction; it must equal $\ln\varphi$ for the theory to be
consistent.  This calibration check is the falsifiable
claim, not the SI numerical value itself.
\end{remark}

\paragraph{Predictions.}
\begin{itemize}
  \item The ratio $k_R/k_B$ (in RS-native units) should match
  $\ln\varphi$ when both are expressed in the same scale.
  \item Entropy in RS is explicitly the cost of information; no
  separate ``entropy constant'' is needed beyond $k_R$.
  \item Thermodynamic irreversibility is $J$-cost minimization:
  the arrow of time points toward lower ledger cost.
\end{itemize}

\section{SI Calibration}

The SI Boltzmann constant $k_B = 1.381 \times 10^{-23}\;\mathrm{J/K}$
is recovered by matching the RS-native thermal energy scale to SI
energy and temperature scales through the calibration seam.

The calibration proceeds in two steps:
\begin{enumerate}
  \item \textbf{Energy anchor}: In RS-native units ($c = 1$,
  $\tau_0 = 1$, $\hbar = \varphi^{-5}$), the coherence energy is
  $E_\mathrm{coh} = \varphi^{-5}$.
  One RS energy unit (``coh'') is calibrated to the Planck energy
  $E_P = \sqrt{\hbar c^5/G} \approx 1.956 \times 10^9\;\mathrm{J}$,
  so $E_\mathrm{coh}^{\mathrm{SI}} = \varphi^{-5}\,E_P$.
  \item \textbf{Temperature anchor}: The RS temperature $T = 1$
  is identified with the Planck temperature
  $T_P = E_P/k_B \approx 1.417 \times 10^{32}\;\mathrm{K}$.
\end{enumerate}
The ratio $k_B = E_P/T_P = E_P/(E_P/k_B)$ is circular by construction
in SI; the RS content is that $k_R = k_B$ in Planck units (i.e.,
$k_B/E_P \cdot T_P = 1$ in Planck units, which holds by definition).
The non-trivial RS prediction is the \emph{structure} $k_R = \ln\varphi$,
which characterizes the fundamental information cost per ledger rung.

\begin{remark}[What RS predicts and does not predict]
RS predicts $k_R = \ln\varphi$ in RS-native units.  It does not
predict the SI numerical value $1.381 \times 10^{-23}\;\mathrm{J/K}$,
which is convention-dependent (determined by the historical definition
of the Kelvin via water's triple point and then fixed by the 2019 SI
revision).  Note that in standard Planck units (where
$c = \hbar = G = k_B = 1$ by convention), $k_B = 1$ by fiat.
RS-native units differ: $c = 1$, $\hbar = \varphi^{-5}$,
$G = \varphi^5$, and $k_R = \ln\varphi \approx 0.481$.  The
falsifiable RS claim is the dimensionless identity
$k_R = \ln\varphi$, which relates the thermal constant to
the $\varphi$-ladder spacing with zero free parameters.
\end{remark}

\section{Conclusion}

The Boltzmann constant is derived from the $\varphi$-ladder: $k_R =
\ln\varphi = J_{\mathrm{bit}}$.  It is the cost of one bit of
information on the recognition ledger.  Temperature is the number of
excited bits; energy is the total bit-cost.

\begin{thebibliography}{9}

\bibitem{Aczel1966}
J.~Acz\'el, \emph{Lectures on Functional Equations and Their Applications},
Academic Press, New York (1966).

\bibitem{CODATA}
E.~Tiesinga, P.~J.~Mohr, D.~B.~Newell, and B.~N.~Taylor,
``CODATA Recommended Values of the Fundamental Physical
Constants: 2022,'' Rev.\ Mod.\ Phys.\ \textbf{96}, 025004 (2024).

\bibitem{Shannon1948}
C.~E.~Shannon, ``A Mathematical Theory of Communication,''
Bell Syst.\ Tech.\ J.\ \textbf{27}, 379 (1948).

\bibitem{RS_Axioms}
J.~Washburn, ``The Algebra of Reality: A Recognition Science Derivation
of Physical Law,'' \emph{Axioms} \textbf{15}(2), 90 (2026).
\texttt{doi:10.3390/axioms15020090}

\end{thebibliography}

\end{document}
